\documentclass[12pt]{article}

\usepackage[utf8]{inputenc}
\usepackage[T1]{fontenc}
\usepackage[spanish]{babel}
\usepackage{lmodern}
\usepackage{booktabs}
\usepackage{amsmath}
\usepackage{forest}
\usepackage{float}
\usepackage{listings}
\usepackage{xcolor}
\usepackage{tikz}
\usepackage{graphicx}
\usepackage{hyperref}

\definecolor{codegreen}{rgb}{0,0.6,0}
\definecolor{codegray}{rgb}{0.5,0.5,0.5}
\definecolor{codepurple}{rgb}{0.58,0,0.82}
\definecolor{backcolour}{rgb}{0.95,0.95,0.92}

\lstdefinestyle{mystyle}{
    backgroundcolor=\color{backcolour},   
    commentstyle=\color{codegreen},
    keywordstyle=\color{magenta},
    numberstyle=\tiny\color{codegray},
    stringstyle=\color{codepurple},
    basicstyle=\ttfamily\footnotesize,
    breakatwhitespace=false,         
    breaklines=true,                 
    captionpos=b,                    
    keepspaces=true,                 
    numbers=left,                    
    numbersep=5pt,                  
    showspaces=false,                
    showstringspaces=false,
    showtabs=false,                  
    tabsize=2
}

\lstset{style=mystyle}

\sloppy
\setlength{\parindent}{0pt}

\begin{document}

\begin{center}
  {\LARGE \textbf{PCA: Principal Component Analysis}}\\[0.5em]
  {Analítica de Datos, Universidad de San Andrés}
\end{center}

\vspace{1em}

Si encuentran algún error en el documento o hay alguna duda, mandenmé un mail a rodriguezf@udesa.edu.ar y lo revisamos.

\section{El Problema de la Alta Dimensionalidad}

Imaginá que estás trabajando con un dataset de vinos que tiene 13 propiedades químicas: acidez fija, acidez volátil, ácido cítrico, azúcar residual, cloruros, dióxido de azufre libre, dióxido de azufre total, densidad, pH, sulfatos, alcohol, etc. O peor: imaginá un dataset con cientos o miles de variables.

\vspace{0.5em}

Esto genera varios problemas:

\begin{itemize}
    \item \textbf{Imposible de visualizar}: no podés graficar 13 dimensiones en una pantalla 2D
    \item \textbf{Multicolinealidad}: muchas de estas propiedades están correlacionadas (por ejemplo, acidez con pH, alcohol con densidad)
    \item \textbf{Costo computacional alto}: entrenar modelos con muchas features es lento y consume mucha memoria
    \item \textbf{Curse of dimensionality}: en espacios de alta dimensión, los datos se vuelven ``dispersos'' y las distancias pierden significado
    \item \textbf{Overfitting}: más features = más riesgo de sobreajuste, especialmente si tenés pocas observaciones
\end{itemize}

La pregunta es: \textbf{¿Podemos reducir el número de features sin perder demasiada información?} Acá entra PCA.

\section{¿Qué es PCA?}

PCA (\textit{Principal Component Analysis} o Análisis de Componentes Principales) es una técnica de \textbf{reducción de dimensionalidad}. Su objetivo es transformar un conjunto de variables (posiblemente correlacionadas) en un conjunto más pequeño de variables no correlacionadas llamadas \textbf{componentes principales}.

\subsection{La Idea Intuitiva}

Pensalo así: tenés un montón de features que están correlacionadas. Por ejemplo, en nuestro dataset de vinos:

\begin{itemize}
    \item Acidez fija, acidez volátil, ácido cítrico, pH $\rightarrow$ todas relacionadas con la ``composición ácida''
    \item Alcohol, densidad, azúcar residual $\rightarrow$ todas relacionadas con la ``composición alcohólica''
    \item Sulfatos, dióxido de azufre libre/total $\rightarrow$ relacionadas con ``conservantes''
\end{itemize}

En lugar de tener todas estas variables, PCA crea nuevas variables (componentes principales) que capturen la mayor parte de la variabilidad:

\begin{itemize}
    \item \textbf{Componente Principal 1 (PC1)}: captura la dirección de máxima variabilidad en los datos (ej: ``balance ácido general del vino'')
    \item \textbf{Componente Principal 2 (PC2)}: captura la segunda dirección de máxima variabilidad, \textbf{ortogonal a PC1} (ej: ``contenido alcohólico y dulzor'')
    \item Y así sucesivamente...
\end{itemize}

La clave es que los primeros componentes principales suelen capturar la mayor parte de la información, entonces podés quedarte solo con ellos y descartar el resto.

\subsection{¿Por Qué Usar PCA?}

Las principales motivaciones para usar PCA son:

\begin{enumerate}
    \item \textbf{Visualización}: reducir a 2 o 3 dimensiones para poder graficar y explorar los datos visualmente
    \item \textbf{Reducir costo computacional}: menos features = entrenamiento más rápido y menos memoria
    \item \textbf{Preprocesamiento para otros modelos}: eliminar multicolinealidad antes de aplicar regresión lineal, por ejemplo
    \item \textbf{Eliminación de ruido}: si las últimas componentes tienen poca varianza, probablemente sean ruido
    \item \textbf{Compresión de datos}: almacenar datos de forma más eficiente
\end{enumerate}

\section{¿Cómo Funciona PCA?}

Vamos a explicar la intuición sin meternos demasiado en el álgebra lineal.

\subsection{Paso a Paso}

\textbf{\underline{Paso 1: Estandarizar los datos}}

\vspace{0.5em}

PCA es sensible a la escala de las variables. Si una variable tiene un rango mucho mayor que otra, va a dominar la varianza. Por eso, primero estandarizamos todas las features para que tengan media 0 y varianza 1.

\[
z_i = \frac{x_i - \mu}{\sigma}
\]

\textbf{\underline{Paso 2: Calcular la matriz de covarianza}}

\vspace{0.5em}

La matriz de covarianza nos dice cómo varían las features entre sí. Si dos features tienen covarianza alta (positiva o negativa), significa que están correlacionadas.

\vspace{0.5em}

\textbf{\underline{Paso 3: Calcular los eigenvectors y eigenvalues}}

\vspace{0.5em}

Acá viene la magia del álgebra lineal. Los \textbf{eigenvectors} (autovectores) nos dan las \textbf{direcciones} de máxima varianza, y los \textbf{eigenvalues} (autovalores) nos dicen \textbf{cuánta varianza} hay en cada dirección.

\begin{itemize}
    \item El eigenvector con el eigenvalue más grande define la primera componente principal (PC1)
    \item El segundo eigenvector define PC2, y así sucesivamente
    \item Los eigenvectors son ortogonales entre sí (perpendiculares), lo que garantiza que los componentes principales no estén correlacionados
\end{itemize}

\textbf{\underline{Paso 4: Proyectar los datos}}

\vspace{0.5em}

Una vez que tenemos los eigenvectors, proyectamos los datos originales sobre estos nuevos ejes. Cada componente principal es una combinación lineal de las features originales:

\[
PC_1 = w_{1,1} \cdot x_1 + w_{1,2} \cdot x_2 + \ldots + w_{1,p} \cdot x_p
\]

donde $w_{1,j}$ son los pesos (componentes del eigenvector correspondiente).

\vspace{0.5em}

\textbf{\underline{Paso 5: Seleccionar los K primeros componentes}}

\vspace{0.5em}

Finalmente, nos quedamos solo con los primeros K componentes principales (los que tienen mayor varianza). Así reducimos de $p$ features originales a $K$ componentes, con $K \ll p$.

\subsection{Ejemplo Simplificado}

Supongamos que tenemos datos de estudiantes con dos features: ``Horas de Estudio'' y ``Nota en el Examen''. Estas variables están muy correlacionadas (quien estudia más, saca mejor nota). Si graficamos los datos en 2D, vamos a ver que los puntos están dispersos a lo largo de una diagonal. PCA encuentra esa diagonal como la primera componente principal (la dirección de máxima varianza). La segunda componente principal sería perpendicular a esa diagonal, capturando pequeñas variaciones (ej: estudiantes que estudian mucho pero les fue mal por nervios, o estudiantes que estudiaron poco pero tuvieron suerte). Si queremos reducir a 1 dimensión, nos quedamos solo con PC1 (la diagonal), que captura la mayor parte de la información.

\section{Varianza Explicada}

Una pregunta clave es: \textbf{¿Con cuántos componentes principales nos quedamos?} Para responder esto, miramos la \textbf{varianza explicada} por cada componente. Recordá que cada eigenvalue nos dice cuánta varianza captura ese componente.

\vspace{0.5em}

La proporción de varianza explicada por el componente $i$ es:

\[
\text{Varianza Explicada}_i = \frac{\lambda_i}{\sum_{j=1}^{p} \lambda_j}
\]

donde $\lambda_i$ es el eigenvalue del componente $i$, y $p$ es el número total de features.

\subsection{Regla del 80-90\%}

Una regla práctica es quedarse con los componentes que en total expliquen al menos el 80-90\% de la varianza total. Por ejemplo:

\begin{itemize}
    \item PC1 explica 50\% de la varianza
    \item PC2 explica 25\%
    \item PC3 explica 15\%
    \item PC4 explica 5\%
    \item PC5+ explican menos del 5\%
\end{itemize}

En este caso, con 3 componentes (PC1, PC2, PC3) ya cubrís el 90\% de la varianza, entonces podés descartar PC4 en adelante.

\subsection{Gráfico de Varianza Acumulada}

Una forma visual de decidir es graficar la \textbf{varianza acumulada}. Buscás el punto donde la curva se ``aplana'', similar al método del codo en K-Means.

\section{Implementación en Python}

Veamos cómo aplicar PCA en Python usando \texttt{scikit-learn}.

\subsection{Imports Necesarios}

\begin{lstlisting}[language=Python]
import numpy as np
import pandas as pd
import matplotlib.pyplot as plt
from sklearn.decomposition import PCA
from sklearn.preprocessing import StandardScaler
\end{lstlisting}

\subsection{Estandarizar y Aplicar PCA}

\begin{lstlisting}[language=Python]
# Supongamos que X es nuestro DataFrame con las features
# Por ejemplo: X = df[['feature1', 'feature2', ..., 'feature64']]

# Paso 1: Estandarizar
scaler = StandardScaler()
X_scaled = scaler.fit_transform(X)

# Paso 2: Aplicar PCA (por ahora, con todos los componentes)
pca = PCA()
X_pca = pca.fit_transform(X_scaled)

# Ver cuanta varianza explica cada componente
print("Varianza explicada por cada componente:")
print(pca.explained_variance_ratio_)

# Varianza acumulada
print("\nVarianza acumulada:")
print(np.cumsum(pca.explained_variance_ratio_))
\end{lstlisting}

\subsection{Reducir a K Componentes}

\begin{lstlisting}[language=Python]
# Decidimos quedarnos con 2 componentes para visualizacion
pca = PCA(n_components=2)
X_pca = pca.fit_transform(X_scaled)

print(f"Forma original: {X_scaled.shape}")
print(f"Forma reducida: {X_pca.shape}")
print(f"Varianza explicada: {pca.explained_variance_ratio_.sum():.2%}")
\end{lstlisting}

\subsection{Visualizar Varianza Explicada}

\begin{lstlisting}[language=Python]
# Grafico de varianza explicada por componente
plt.figure(figsize=(10, 6))
plt.bar(range(1, len(pca.explained_variance_ratio_) + 1),
        pca.explained_variance_ratio_)
plt.xlabel('Componente Principal')
plt.ylabel('Varianza Explicada')
plt.title('Varianza Explicada por Componente')
plt.show()

# Grafico de varianza acumulada
plt.figure(figsize=(10, 6))
plt.plot(range(1, len(pca.explained_variance_ratio_) + 1),
         np.cumsum(pca.explained_variance_ratio_), marker='o')
plt.xlabel('Numero de Componentes')
plt.ylabel('Varianza Acumulada')
plt.title('Varianza Acumulada')
plt.axhline(y=0.9, color='r', linestyle='--', label='90% varianza')
plt.legend()
plt.grid(True)
plt.show()
\end{lstlisting}

\section{Interpretación de Componentes Principales}

Los componentes principales son combinaciones lineales de las features originales. Podés ver qué features contribuyen más a cada componente mirando los ``loadings'' (pesos).

\begin{lstlisting}[language=Python]
# Ver los loadings del primer componente principal
loadings = pca.components_[0]
feature_names = X.columns  # Asumiendo que X es un DataFrame

# Crear DataFrame de loadings
loadings_df = pd.DataFrame({
    'feature': feature_names,
    'loading': loadings
}).sort_values('loading', key=abs, ascending=False)

print("Contribucion de cada feature a PC1:")
print(loadings_df.head(10))
\end{lstlisting}

Por ejemplo, si estás trabajando con datos de vinos y ves que PC1 tiene loadings altos en ``acidez'', ``pH'', y ``contenido alcohólico'', podés interpretar PC1 como un eje de ``composición química''.

\section{Reconstrucción de Datos}

Una vez que aplicaste PCA y te quedaste con K componentes, podés \textbf{reconstruir} una aproximación de los datos originales. Obviamente, si K < número de features originales, va a haber pérdida de información.

\begin{lstlisting}[language=Python]
# Reconstruir datos originales desde componentes principales
X_reconstructed = pca.inverse_transform(X_pca)

# Calcular error de reconstruccion (MSE)
mse = np.mean((X_scaled - X_reconstructed) ** 2)
print(f"Error de reconstruccion (MSE): {mse:.4f}")
\end{lstlisting}

La reconstrucción es útil para compresión de datos (ej: compresión de imágenes).

\section{PCA en el Pipeline de Machine Learning}

PCA se puede usar como preprocesamiento antes de entrenar un modelo supervisado. Esto puede:

\begin{itemize}
    \item Reducir el tiempo de entrenamiento
    \item Mitigar overfitting (menos features = menos complejidad)
    \item Eliminar multicolinealidad
\end{itemize}

\subsection{Ejemplo de Pipeline}

\begin{lstlisting}[language=Python]
from sklearn.model_selection import train_test_split
from sklearn.linear_model import LogisticRegression
from sklearn.pipeline import Pipeline

# Dividir datos
X_train, X_test, y_train, y_test = train_test_split(
    X, y, test_size=0.2, random_state=42
)

# Crear pipeline: Scaler -> PCA -> Modelo
pipeline = Pipeline([
    ('scaler', StandardScaler()),
    ('pca', PCA(n_components=10)),
    ('classifier', LogisticRegression(max_iter=1000))
])

# Entrenar
pipeline.fit(X_train, y_train)

# Evaluar
score = pipeline.score(X_test, y_test)
print(f"Accuracy con PCA: {score:.3f}")
\end{lstlisting}

Podés comparar el performance con y sin PCA para ver si vale la pena.

\section{Ventajas y Limitaciones de PCA}

\subsection{Ventajas}

\begin{itemize}
    \item \textbf{Reduce dimensionalidad} sin necesidad de seleccionar manualmente features
    \item \textbf{Elimina multicolinealidad}: los componentes principales son ortogonales (no correlacionados)
    \item \textbf{Visualización}: permite graficar datos de alta dimensión en 2D o 3D
    \item \textbf{Eficiencia computacional}: reduce el costo de entrenar modelos
    \item \textbf{Elimina ruido}: componentes con poca varianza suelen ser ruido
\end{itemize}

\subsection{Limitaciones}

\begin{itemize}
    \item \textbf{Pérdida de interpretabilidad}: los componentes principales son combinaciones de features, difíciles de interpretar
    \item \textbf{Asume linealidad}: PCA solo encuentra relaciones lineales. Si las features están relacionadas de forma no lineal, PCA no va a funcionar bien (para eso existe kernel PCA o t-SNE)
    \item \textbf{Sensible a outliers}: los componentes pueden estar distorsionados por valores extremos
    \item \textbf{Requiere estandarización}: sin estandarizar, las features con mayor rango dominan
    \item \textbf{No supervisa}: PCA no considera la variable objetivo ($y$). Podría eliminar features importantes para la predicción
\end{itemize}

\section{¿Cuándo Usar PCA?}

\subsection{Casos de Uso Ideales}

Usá PCA cuando:

\begin{itemize}
    \item Tenés muchas features correlacionadas
    \item Querés visualizar datos de alta dimensión
    \item El entrenamiento de modelos es muy lento por la cantidad de features
    \item Tenés problemas de multicolinealidad (ej: antes de regresión lineal)
    \item Querés hacer compresión de datos (ej: imágenes, audio)
\end{itemize}

\subsection{Casos Donde NO Usar PCA}

No uses PCA si:

\begin{itemize}
    \item Tus features no están correlacionadas (PCA no va a reducir mucho)
    \item La interpretabilidad es crítica (los componentes son difíciles de interpretar)
    \item Las relaciones entre features son no lineales (considerá kernel PCA o autoencoders)
    \item Tenés muy pocas features (no tiene sentido reducir si ya tenés 3-4 variables)
\end{itemize}

\section{PCA vs Otras Técnicas de Reducción de Dimensionalidad}

Existen otras técnicas además de PCA:

\begin{itemize}
    \item \textbf{t-SNE}: mejor para visualización, pero no para reducción computacional (es lento)
    \item \textbf{UMAP}: similar a t-SNE pero más rápido y escalable
    \item \textbf{Feature Selection}: seleccionar features importantes en lugar de crear combinaciones (más interpretable)
    \item \textbf{Autoencoders}: redes neuronales que aprenden representaciones comprimidas (no lineales)
\end{itemize}

PCA es generalmente la primera opción por su simplicidad y velocidad.

\section{Ejemplo Aplicado: PCA en Imágenes}

Un caso de uso clásico de PCA es compresión de imágenes. Cada píxel es una feature, y PCA encuentra patrones en cómo varían los píxeles.

\subsection{Dataset de Dígitos}

El dataset \texttt{digits} de sklearn tiene imágenes de 8x8 píxeles (64 features). Con PCA, podés reducir a 10-20 componentes sin perder mucha información.

\begin{lstlisting}[language=Python]
from sklearn.datasets import load_digits

# Cargar datos
digits = load_digits()
X = digits.data  # 1797 imagenes x 64 pixels
y = digits.target  # Digitos 0-9

# Aplicar PCA
pca = PCA(n_components=20)
X_pca = pca.fit_transform(X)

print(f"Forma original: {X.shape}")
print(f"Forma reducida: {X_pca.shape}")
print(f"Varianza explicada: {pca.explained_variance_ratio_.sum():.2%}")
\end{lstlisting}

Con solo 20 componentes (de 64 originales), podés capturar más del 90\% de la varianza, lo que significa que las imágenes reconstruidas van a ser muy parecidas a las originales.

\section{Recursos Extra}

Si querés profundizar más en PCA, te recomiendo:

\begin{itemize}
    \item \textbf{StatQuest: PCA clearly explained}: \url{https://www.youtube.com/watch?v=FgakZw6K1QQ}
    \item \textbf{3Blue1Brown: Eigenvectors and eigenvalues}: \url{https://www.youtube.com/watch?v=PFDu9oVAE-g}
    \item \textbf{Documentación sklearn}: \url{https://scikit-learn.org/stable/modules/decomposition.html#pca}
\end{itemize}

\end{document}
